\documentclass[10pt,a4paper]{amsart}
\usepackage[margin=19mm]{geometry}
\usepackage{amsmath,amssymb,mathtools}
\usepackage[scr=rsfs,cal=euler]{mathalfa}
\usepackage{xfrac}
\usepackage[unicode,hidelinks,bookmarks=false]{hyperref}
\newcommand{\ie}{i.e.\ }
\newcommand{\cf}{cf.\ }
\newcommand{\resp}{respectively\ }
\newcommand{\Subsection}{\subsection}
\providecommand{\nobreakdash}{\nobreak\hbox{-}}
\setlength{\emergencystretch}{2em}
\multlinegap0pt
\usepackage{tikz}
\usetikzlibrary{cd}
\usepackage{amssymb}
\newtheorem{lemma}{Lemma}
\title{Gabriel's Theorem for Infinite Quivers}
\author{Nathaniel Gallup, Stephen Sawin}
\date{Source: arXiv:2312.09904v1 (CC BY 4.0)}
\begin{document}
\maketitle

\noindent\emph{Source and licence.} Gabriel's Theorem for Infinite Quivers. Version arXiv:2312.09904v1 (CC BY 4.0), distributed by arXiv under the Creative Commons Attribution 4.0 International licence, http://creativecommons.org/licenses/by/4.0/, as recorded in the arXiv licence metadata for this exact version and shown on the abstract page https://arxiv.org/abs/2312.09904v1. Full licence notice: \texttt{licenses/arxiv-2312.09904-CC-BY-4.0.txt}.\par\smallskip

\noindent\textit{Editorial scope.} Counted unit: the article's lemma types (source0.tex lines 396--398). The article's complete original proof of it is delivered with it (source0.tex lines 400--453, one \texttt{proof} environment of 54 lines). No mathematics is written, repaired or re-derived: the statement and the proof are the article's own lines, in the article's own notation, and no retained line is substituted or re-flowed.  No further source-local context is needed: every cross-reference of every retained line resolves inside the two delivered spans.  Nothing was omitted from any retained span, so the package carries no omission record.  No retained line contains a citation, so the wrapper carries no bibliography and no bibliography entry.  The retained lines contain the article's own diagram code, which the wrapper typesets with the standard tikz packages; no image file is delivered and the package declares no asset dependency.  Editorial formatting only, in addition: an AMS \texttt{amsart} preamble that restates the article's own declarations and macros used by the retained lines, namely the environments \texttt{lemma} on separate counters, together with the standard packages those lines need.  The retained environments print in this document as Lemma 1 in order of appearance, whereas the article prints the same items with the numbering of its own full document.  The complete native source of the article as distributed in the arXiv e-print for this exact version is bundled unchanged as \texttt{sources/151564/source0.tex}.
\begin{lemma}\label{lm: ordering among types}
    If there is a nontrivial map from $\alpha$ to $\beta$ which is either injective or surjective, then $\alpha$ has type less than or equal to $\beta$. Therefore in an infinite chain of nontrivial maps between indecomposables, eventually all the indecomposables are of the same type.  
\end{lemma}

\begin{proof}
We consider the different possible types of $\beta$ as different cases. 

Case 1: Suppose that $\beta$ is of Type I. We claim that $\alpha$ must also be of Type I. If not then $\alpha_k \neq 0$. Since there is a nonzero map $\alpha \to \beta$, the supports of $\alpha$ and $\beta$ must intersect, therefore there exists some vertex $\ell$ where $\alpha_\ell \neq 0$ and $\beta_\ell \neq 0$. But the support of $\alpha$ is connected, so there is a directed path from $k$ to $\ell$ contained in the support of $\alpha$. Because $\beta_k = 0$ there must be an arrow $i \to j$ along this path with $\alpha_i , \alpha_j \neq 0$, $\beta_i=0$ and $\beta_j=1$. Hence we obtain the following diagram for representatives $V$ and $W$ of $\alpha$ and $\beta$ respectively.

\begin{equation}\label{eq: first diagram}
\begin{tikzcd}
V_j \neq 0 \arrow[d]  & V_i \neq 0 \arrow[l]  \arrow[d] & 
\\ W_j \cong \mathbb{F} & W_i = 0 \arrow[l] 
\end{tikzcd}
\end{equation}

Since the morphism $V \to W$ is either injective or surjective and $\beta_i = 0$, the map $V_i \to W_i$ must in fact be surjective. Since the map $V_i \to V_j$ is nonzero, this diagram cannot commute, contradiction. 

Case 2: Suppose that $\beta$ is of Type II. We claim that $\alpha$ cannot have $\alpha_k = 2$ and also cannot have $\alpha_j = 0$ for any $j$ adjacent to $k$ (this will show that $\alpha$ cannot be of Types III (a), III(b), IV, or V). The latter creates the following diagram

\begin{equation}\label{eq: one to zero bottom iso}
\begin{tikzcd}
V_j = 0 \arrow[d]  & V_i \cong \mathbb{F} \arrow[l]  \arrow[d] & 
\\ W_j \neq 0 & W_i \neq 0 \arrow[l, "\sim" '] 
\end{tikzcd}
\end{equation}

with $k = i$ which is impossible because the down arrows are both injections and the bottom arrow is an isomorphism hence $V_i \to W_i \to W_j$ is nonzero. In the former case, if $i$ and $j$ are the two vertices on the short arms, we obtain a diagram of the form 

\begin{equation}\label{eq: second diagram}
\begin{tikzcd}
V_i \cong \mathbb{F} \arrow[d]  & V_k \cong \mathbb{F}^2 \arrow[l]  \arrow[d] \arrow[r] & V_j \cong \mathbb{F} \arrow[d]  
\\ W_i \cong \mathbb{F} & W_k \cong \mathbb{F} \arrow[l] \arrow[r]  & W_j \cong \mathbb{F}
\end{tikzcd}
\end{equation}

where all maps are surjective and the outside and bottom maps are isomorphisms. Therefore the kernel of the map $V_k \to W_k$ must be equal to the kernel of the map $V_k \to V_j$ and to the kernel of the map $V_k \to V_i$. But these kernels are transverse lines in $\mathbb{F}$ and in particular are not equal so we obtain a contradiction. 

Case 3: Suppose that $\beta$ is of Type III (a), i.e. that $\beta_k=1$ and $\beta_j=1$ for exactly two $j$'s adjacent to $k$. If $\alpha$ has Type IV or V then this again creates Diagram \ref{eq: first diagram} with $V_i \to W_j \to W_j$ nonzero which is impossible.

Case 4: Now suppose that $\beta$ is of Type III (b), i.e. that $\beta_k = 2$. If $\alpha$ is of Type IV or V then there are two vertices $i$ and $j$ adjacent to $k$ with $a_i = a_j = 0$. Hence we obtain a diagram of the following form. 

\begin{equation}\label{eq: two zeros mapping to thick}
\begin{tikzcd}
V_i = 0 \arrow[d]  & V_k \cong \mathbb{F} \arrow[l]  \arrow[d] \arrow[r] & V_j = 0 \arrow[d]  
\\ W_i \cong \mathbb{F} & W_k \cong \mathbb{F}^2 \arrow[l] \arrow[r]  & W_j \cong \mathbb{F}
\end{tikzcd}
\end{equation}


Hence the image of the map $V_k \to W_k$ must be equal both to the kernel of the map $W_k \to W_i$ and to that of $W_k \to W_j$. Again these kernels are in fact transverse lines so this is impossible.  

Case 5: Suppose that $\beta$ is of Type IV. If $\alpha$ is of Type V, this again creates Diagram \ref{eq: first diagram} with $V_i \to W_j \to W_j$ nonzero which is impossible.

%If $\beta$ is of type I, then there is an arrow $i \to j$ with $\beta_j=1$ and $\beta_i=0$.  If $\alpha$ is type II or higher, then $\alpha_i>0$, and the map from $i$ to $j$ is surjective in $\alpha$.  Therefore the  map from $\alpha_j$ to $\beta_j$ must be trivial, so the map from $\alpha$ to $\beta$ must be trivial.  If $\alpha$ is of type IV, then $\alpha_k=1$ and  $\alpha_j=1$ for all $j$ adjacent to $k$. If $\alpha$ and $\beta$ are both of type II,III,V or VI, that is if $k$ has dimension $1$, then a nontrivial map between them must be an isomorphism on $k$, and therefore if $\alpha_j=1$ then $\beta_j=1$ for any adjacent $k$, and thus the type of $\alpha$ must be less than the type of $\beta$.  If $\beta$ is of type IV and $\alpha$ is neither I nor IV, then the map between $\alpha$ and $\beta$ is an injection, and because in the two dimensional space $\beta$ assigns to $k$ the kernels of the three arrow maps are transverse lines, any vector in that space is in the kernel of at most one of the maps, and therefore $\alpha$ is of type II or III.  Similarly if $\alpha$ is of type IV, and $\beta$ has $\beta_k=1$, then the map between $\alpha$ and $\beta$ is an surjection,  and because in the two dimensional space $\beta$ assigns to $k$ the kernels of the three arrow maps are transverse lines, any one-dimensional quotient identifies the remaining one-dimensional space with at least two of the kernels, and therefore $\beta$ is of type V or VI.

  
\end{proof}

\end{document}
